%*******************************************************************************
%****************************** Chapter 7 *************************************
%*******************************************************************************

\chapter{System Integration, Validation, and Comparative Analysis}
\label{ch:validation}

\ifpdf
    \graphicspath{{Chapter7/Figs/Raster/}{Chapter7/Figs/PDF/}{Chapter7/Figs/}}
\else
    \graphicspath{{Chapter7/Figs/Vector/}{Chapter7/Figs/}}
\fi

This chapter reports the integrated behaviour of MIXR-1, compares it against the legacy Armfield CEK-A on the parameters defined in Section~\ref{subsec:scope}, quantifies measurement error against the independent references of Section~\ref{sec:arch-validation}, and examines the sources of residual error.


%********************************** 7.1 **************************************
\section{Integrated Test Overview}
\label{sec:val-overview}

% TODO: Table 7.1 -- master results table: every test case (TC-U, TC-S, TC-H), its pass criterion,
% the measured result, and the section reporting it. This is the single table a panel will look
% for, so it must be complete and must agree with the subsystem chapters.


%********************************** 7.2 **************************************
\section{Comparative Analysis: MIXR-1 versus the Armfield CEK-A}
\label{sec:val-comparison}

\subsection{Basis of Comparison}
\label{subsec:val-basis}

% TODO: state plainly how the CEK-A figures were obtained: direct measurement on the unit, its
% documented specifications, the laboratory manual, or historical student data. If the unit could
% not be operated for any parameter, say so and name the substitute source. The credibility of
% this whole section rests on that disclosure.

\subsection{Benchmarking Against Literature}
\label{subsec:val-literature}

% TODO: compare the MIXR-1 N_Re-N_Po curve against the Rushton correlation as tabulated by Brown
% et al., and against the CEK-A data for the same experiment. Figure 7.1 -- all three on one axis.
% Carry forward the geometric-mismatch caveat from Section~\ref{subsec:agit-plausibility}.

\subsection{Parameter-by-Parameter Comparison}
\label{subsec:val-parameters}

% TODO: Table 7.2 -- one row per parameter from the scope: speed-setpoint tracking, steady-state
% speed error, torque measurement method and error, power computation, level sensing, data-logging
% completeness, and MATLAB/Simulink integration. Columns: CEK-A, MIXR-1, and the basis of each
% figure. Keep the claims to what the tests in Chapters 5 and 6 actually measured.

The most substantive differences are structural rather than numerical. The CEK-A sets speed open loop, so a target speed is reached by manual search and verified against an analog tachometer that the laboratory manual corrects with an empirical expression, as described in Section~\ref{subsec:che3214l}. MIXR-1 closes the loop on a digital measurement, so the setpoint is commanded directly. The CEK-A estimates torque from a spring dynamometer reading and the torque-arm geometry; MIXR-1 measures reaction torque directly, which is the method \citet{ascanio2004} identify as appropriate at this scale. And the CEK-A records all data manually, whereas MIXR-1 logs continuously and exports in machine-readable form.

\subsection{Cost and Software Openness}
\label{sec:val-cost}

% TODO: Table 7.3 -- cost comparison: CEK-A at its CEPCI-adjusted acquisition cost, the CEK-MKII
% quotation, and the MIXR-1 per-unit materials cost from Appendix~\ref{app:bom}. Derive the
% adjustment in Appendix~\ref{app:calculations} and resolve the base-year question flagged in
% Chapter 1. State how many MIXR-1 units a CEK-MKII budget would fund.

Beyond cost, the two differ in what they expose. The commercial unit's bundled software logs and plots but does not expose raw sensor data or support deployment of a student-authored controller. MIXR-1 streams raw records over a documented interface and accepts deployed Simulink models, which is what makes the single apparatus serve both target courses.


%********************************** 7.3 **************************************
\section{Percentage-Error Analysis}
\label{sec:val-error}

% TODO: Table 7.4 -- consolidated error summary: measured quantity, reference, mean percent error,
% maximum percent error, and the range over which it was evaluated. Gather from TC-U1 through
% TC-U7 rather than recomputing.

% TODO: propagate the speed and torque errors into the uncertainty in P and N_Po, continuing the
% treatment begun in Section~\ref{subsec:agit-power}, and state which term dominates.


%********************************** 7.4 **************************************
\section{Error Sources and Mitigation}
\label{sec:val-errorsources}

% TODO: Table 7.5 -- error source, affected quantity, estimated magnitude, mitigation applied, and
% residual effect. Candidate sources identified during development:
%   - timestamp quantization and scheduling jitter in the speed estimate (Sections 4.6, 5.2)
%   - encoder quadrature non-idealities, largely cancelled by slot alignment (Section 5.2)
%   - torque amplifier zero drift, reduced by automatic tare (Section 5.5)
%   - uncertainty in the theoretical torque reference itself (Section 5.5)
%   - ToF return from a disturbed or specular water surface (Section 6.2)
%   - geometric departure from the standard correlation configuration (Section 5.8)
%   - fluid property assumptions (temperature-dependent density and viscosity) in N_Re


%********************************** 7.5 **************************************
\section{Discussion}
\label{sec:val-discussion}

% TODO: answer the objectives of Section~\ref{sec:objectives} one at a time, with the numbers.
% Then state honestly what the apparatus does not yet do well, and what the limits of this
% validation are: internal testing only, no classroom deployment, no multi-semester reliability
% data. Those are already declared as limitations, and acknowledging them here is stronger than
% leaving the panel to raise them.